\section{Metodología}

\subsection{Modelos evaluados}

\begin{table}[htbp]
\centering
\caption{Conjunto de modelos evaluados, con la versión de \texttt{transformers} resuelta para cada uno. La línea horizontal separa la franja sub-1B (arriba) de la franja 1--2B (abajo). Forzaron el pin: tokenizer no soportado en 4.57.6 (\texttt{LFM2.5-230M}); íd. (mismo tokenizer) (\texttt{LFM2.5-350M}); regresión de caché en la versión 5.14.1 (\texttt{granite-4.0-350m}); arquitectura no reconocida en 4.57.6 (\texttt{Qwen3.5-0.8B}).}
\label{tab:modelos}
{\footnotesize
\begin{tabular}{lrlr}
\toprule
Modelo & Parámetros & Familia & \texttt{transformers} \\
\midrule
LFM2.5-230M & 0.23B & LiquidAI & 5.14.1 \\
granite-4.0-h-350m & 0.34B & IBM Granite & 4.57.6 \\
LFM2.5-350M & 0.35B & LiquidAI & 5.14.1 \\
granite-4.0-350m & 0.35B & IBM Granite & 4.57.6 \\
SmolLM2-360M-Instruct & 0.36B & HuggingFaceTB & 4.57.6 \\
Qwen3.5-0.8B & 0.80B & Qwen & 5.14.1 \\
\midrule
OLMo-2-0425-1B-Instruct & 1.00B & AllenAI & 4.57.6 \\
LFM2.5-1.2B-Instruct & 1.20B & LiquidAI & 4.57.6 \\
granite-4.0-h-1b & 1.50B & IBM Granite & 4.57.6 \\
Qwen2.5-1.5B-Instruct & 1.54B & Qwen & 4.57.6 \\
granite-4.0-1b & 1.60B & IBM Granite & 4.57.6 \\
SmolLM2-1.7B-Instruct & 1.71B & HuggingFaceTB & 4.57.6 \\
\bottomrule
\end{tabular}
}
\end{table}


Se evaluaron doce modelos abiertos, públicos en HuggingFace y sin
restricción de licencia ni de token, en dos franjas de seis: sub-1B y de 1
a 2B (Tabla
\ref{tab:modelos}). \texttt{Qwen2.5-1.5B-Instruct} (1,54\,B) se incorporó
deliberadamente como representante de la generación anterior de SLMs bajo el
mismo protocolo (Sección \ref{sec:resultados-globales}).

\subsection{Dataset de comandos}

El dataset consta de 32 comandos de domótica en español rioplatense (voseo),
con una anotación de referencia de cinco campos: \texttt{intent},
\texttt{dispositivo}, \texttt{ubicacion}, \texttt{valor} y \texttt{unidad}.
Cada comando pertenece además a una de seis categorías lingüísticas cerradas
— encendido/apagado simple, ajuste con valor numérico explícito, ajuste
relativo o cualitativo, múltiples dispositivos, consulta de estado y negación
o contexto conversacional —, que asigna el juez automático (Sección
\ref{sec:dos-etapas}) y no un anotador humano.

\subsection{Protocolo de prompting y esquema de salida}
\label{sec:prompting}

El prompt de sistema es único y fijo para los doce modelos, aplicado
mediante la plantilla de conversación del propio \emph{tokenizer} de cada
uno. Se transcribe íntegro junto con el prompt de usuario; ambos provienen
de constantes del código, no de una transcripción manual.

\paragraph{Prompt de sistema (texto completo).}
{\footnotesize\ttfamily
\begin{verbatim}
Sos un módulo de interpretación de comandos de voz para un sistema
domótico en español rioplatense. Convertí el comando del usuario en UN
objeto JSON con estos campos: "intent"
(encender/apagar/ajustar/consultar), "dispositivo" (luz,
aire_acondicionado, calefaccion, persiana, cortina, tv, parlante,
cerradura, ventilador, todos), "ubicacion" (living, cocina, dormitorio,
bano, garage, patio, entrada, todas), "valor" (número si el comando lo
da, si no null), "unidad" ("%" o "°C" si corresponde, si no null).
Reglas: respondé SOLO el JSON, sin texto adicional; si el comando ajusta
sin dar un número, "valor" y "unidad" deben ser null.

Ejemplo: Comando: "Poné el aire de la cocina en 20 grados."
{"intent": "ajustar", "dispositivo": "aire_acondicionado", "ubicacion":
"cocina", "valor": 20, "unidad": "°C"}

IMPORTANTE: los campos "intent", "dispositivo", "ubicacion" y "unidad"
son campos CATEGORICOS CERRADOS. Los valores listados entre parentesis
mas arriba son, textualmente, los unicos outputs aceptados para esos
campos: copialos exactamente como aparecen, sin tildes, sin mayusculas,
sin plural y sin reformular. Usar un sinonimo, una traduccion o
cualquier variante que no este literalmente en esa lista (por ejemplo
"televisor" o "tele" en lugar de "tv") es una violacion del formato y la
respuesta se considera incorrecta. Si ningun valor de la lista aplica,
usa null.

Ejemplo 2: Comando: "Bajale un poco a la luz del living."
{"intent": "ajustar", "dispositivo": "luz", "ubicacion": "living",
"valor": null, "unidad": null}
\end{verbatim}
}

\paragraph{Prompt de usuario (uno por comando, tras el de sistema).}
{\footnotesize\ttfamily
\begin{verbatim}
Comando: "<comando del usuario>"
\end{verbatim}
}


El prompt solicita un único objeto JSON con los cinco campos, incluye dos
ejemplos \emph{few-shot} y agrega una cláusula taxativa sobre los campos
categóricos cerrados: los valores listados entre paréntesis son,
textualmente, las únicas salidas aceptadas, y un sinónimo válido en español
—``tele'' por ``tv''— constituye una violación de formato. Los resultados de
la Sección \ref{sec:resultados} se calculan bajo este prompt, por lo que la comparación entre
arquitecturas nunca mezcla prompts.

\subsection{Métricas}
\label{sec:metricas}

La \textbf{exactitud estricta} exige que los cinco campos coincidan
textualmente con la referencia bajo ese protocolo taxativo; la
\textbf{exactitud laxa (semántica)} incorpora además los casos que el juez
etiquetó como \texttt{sin\_error\_semantico} o \texttt{uso\_de\_sinonimos}. La
estricta es la métrica primaria; la laxa es una cota superior que separa
errores de forma de errores de contenido, y la brecha entre ambas es un
resultado del estudio, no una corrección que deba promediarse. Se reportan
además la tasa de JSON válido, la exactitud por campo y por categoría, y la
latencia promedio.

\subsection{Verificación automática en dos etapas}
\label{sec:dos-etapas}

En la \textbf{etapa 1} cada respuesta se compara campo a campo, de forma
determinista, contra la referencia, sin ningún modelo adicional, y de ahí
sale la exactitud estricta. El juez de la \textbf{etapa 2} se determina por
una regla fijada de antemano: el modelo con mayor exactitud estricta en la
etapa 1, con empates resueltos por mayor cantidad de parámetros y luego por
el orden del listado. Ambas etapas usan decodificación \emph{greedy}
(\texttt{do\_sample=False}, temperatura 0); de eso — no de fijar el
identificador del juez de antemano — depende la reproducibilidad.

El juez clasifica cada respuesta con \texttt{match\_exact = False} en un
subconjunto no vacío de siete etiquetas cerradas — confusión de intención, de
dispositivo y de ubicación; alucinación de valor o unidad; valor numérico
incorrecto; más \texttt{sin\_error\_semantico} y \texttt{uso\_de\_sinonimos},
que definen la exactitud laxa — y asigna a cada comando exactamente una
categoría lingüística. Una salida que no valide contra ese vocabulario se
reintenta una vez y, si vuelve a fallar, recibe una etiqueta de respaldo
determinista; su frecuencia se reporta. Al ser uno de los modelos evaluados,
el juez adjudica también sobre sus propias respuestas (Sección \ref{sec:amenazas}).

\subsection{Entorno de ejecución y reproducibilidad}
\label{sec:entorno-repro}

El host es un equipo con procesador Intel Core i5-11400H (6 núcleos físicos / 12 hilos), sin GPU:
los doce modelos y el juez se ejecutan en CPU, cada uno en su propia imagen
de contenedor y de forma secuencial.

El conjunto evaluado necesita \textbf{dos grupos de versión} de
\texttt{transformers}, mutuamente excluyentes (Tabla \ref{tab:modelos}):
4.57.6 para nueve modelos y 5.14.1 para los otros tres.
\texttt{granite-4.0-350m} falla bajo 5.14.1, y \texttt{LFM2.5-230M},
\texttt{LFM2.5-350M} y \texttt{Qwen3.5-0.8B} fallan bajo 4.57.6: el segundo
grupo se resuelve precisamente a la versión que hace fallar al modelo que
impone el primero, de modo que ninguna versión mayor única cubre el conjunto
(Sección \ref{sec:amenazas}).

La configuración de Docker de esta máquina expone únicamente \textbf{dos
procesadores lógicos} y un límite de \textbf{8\,GB}, y las trece ejecuciones
— doce del barrido más la del juez — fijan \texttt{-{}-memory=8g -{}-cpus=2
-{}-cpuset-cpus=0-1}, el mismo par de identificadores lógicos en todas. Que
ese par corresponda a un mismo núcleo físico por \emph{hyperthreading} (SMT)
o a dos núcleos distintos lo decide el hipervisor y no es observable desde el
contenedor. La ejecución es secuencial porque la inferencia en CPU está
limitada por el ancho de banda de memoria, que no se particiona por
contenedor. Ninguno de estos factores afecta a la exactitud, calculada con decodificación
determinista, pero sí a las latencias: comparables entre los doce modelos,
aunque sus valores absolutos no sean los mejores alcanzables en
CPU\footnote{Código y datos: \url{https://github.com/carlosbinker/cacic-2026}
(rama \texttt{feat/reescritura-experimento-2026}).} (Sección \ref{sec:amenazas}).
